% Part: computability
% Chapter: recursive-functions
% Section: pr-functions

\documentclass[../../../include/open-logic-section]{subfiles}

\begin{document}

\olfileid{cmp}{rec}{prf}
\olsection{توابع بازگشتی اولیه}

دوباره ثبت کنیم که چگونه می‌توان با بازگشت اولیه و ترکیب از توابع
موجود، توابع تازه‌ای تعریف کرد.

\begin{defn}
  \ollabel{defn:primitive-recursion} فرض کنید $f$ تابعی $k$موضعی
  ($k\ge 1$) و $g$ تابعی $(k+2)$موضعی باشد. تابعی که با
  \emph{بازگشت اولیه از $f$ و~$g$} تعریف می‌شود، تابع
  $(k+1)$موضعیِ~$h$ است که معادله‌های زیر آن را تعریف می‌کنند:
  \begin{align*}
    h(x_0,\dots,x_{k-1},0) & =  f(x_0,\dots,x_{k-1}) \\
    h(x_0,\dots,x_{k-1},y+1) & = g(x_0,\dots,x_{k-1}, y, h(x_0,\dots,x_{k-1}, y))
  \end{align*}
\end{defn}

\begin{defn}
  \ollabel{defn:composition} فرض کنید $f$ تابعی $k$موضعی و
  $g_0$، \dots،~$g_{k-1}$ تعداد $k$ تابع باشند که همگی $n$موضعی‌اند.
  تابعی که با \emph{ترکیب از $f$ و
  $g_0$، \dots،~$g_{k-1}$} تعریف می‌شود، تابع $n$موضعیِ~$h$ است که
  چنین تعریف می‌شود:
  \[
  h(x_0, \dots, x_{n-1}) =
  f(g_0(x_0, \dots, x_{n-1}), \dots, g_{k-1}(x_0, \dots, x_{n-1})).
  \]
\end{defn}

علاوه بر~$\Succ$ و توابع افکنش
\[
\Proj{n}{i}(x_0,\dots,x_{n-1}) = x_i,
\]
به ازای هر عدد طبیعی $n$ و هر $i<n$، تابع~$\Zero(x)=0$ را نیز در
میان توابع بازگشتی اولیه می‌گنجانیم.

\begin{defn}
  مجموعهٔ توابع بازگشتی اولیه، مجموعهٔ تابع‌های از $\Nat^n$ به
  $\Nat$ است که به‌طور استقرایی با بندهای زیر تعریف می‌شود:
  \begin{enumerate}
  \item $\Zero$ بازگشتی اولیه است.
  \item $\Succ$ بازگشتی اولیه است.
  \item هر تابع افکنش $\Proj{n}{i}$ بازگشتی اولیه است.
  \item اگر $f$ تابعی $k$موضعی و بازگشتی اولیه و
    $g_0$، \dots،~$g_{k-1}$ توابعی $n$موضعی و بازگشتی اولیه باشند،
    آنگاه ترکیب $f$ با $g_0$، \dots،~$g_{k-1}$ بازگشتی اولیه است.
  \item اگر $f$ تابعی $k$موضعی و بازگشتی اولیه و $g$ تابعی
    $(k+2)$موضعی و بازگشتی اولیه باشد، آنگاه تابع تعریف‌شده با بازگشت
    اولیه از $f$ و~$g$ بازگشتی اولیه است.
\end{enumerate}
\end{defn}

\begin{explain}
خلاصه‌تر اینکه مجموعهٔ توابع بازگشتی اولیه کوچک‌ترین مجموعه‌ای است که
$\Zero$، $\Succ$ و توابع افکنش~$\Proj{n}{j}$ را دربر می‌گیرد و نسبت
به ترکیب و بازگشت اولیه بسته است.

راه دیگر توصیف مجموعهٔ توابع بازگشتی اولیه، تعریف آن برحسب «مرحله‌ها»
است. بگذارید $S_0$ مجموعهٔ توابع آغازین، یعنی $\Zero$، $\Succ$ و
افکنش‌ها باشد. این‌ها توابع بازگشتی اولیهٔ مرحلهٔ~$0$ هستند. هرگاه
مرحلهٔ $S_i$ تعریف شده باشد، بگذارید $S_{i+1}$ مجموعهٔ همهٔ توابعی
باشد که با یک بار کاربرد ترکیب یا بازگشت اولیه بر توابع موجود
در~$S_i$ به دست می‌آیند. آنگاه
\[
S = \bigcup_{i \in \Nat} S_i
\]
مجموعهٔ همهٔ توابع بازگشتی اولیه است.
\end{explain}

بررسی کنیم که $\Add$ تابعی بازگشتی اولیه است.

\begin{prop}
  تابع جمع $\Add(x,y)=x+y$ بازگشتی اولیه است.
\end{prop}

\begin{proof}
از پیش تعریفی با بازگشت اولیه برای $\Add$ برحسب دو تابع $f$ و~$g$
داریم که با قالب \olref{defn:primitive-recursion} سازگار است:
\begin{align*}
  \Add(x_0, 0) & = f(x_0) = x_0\\
  \Add(x_0, y+1) & = g(x_0, y, \Add(x_0, y)) =
  \Succ(\Add(x_0, y))
\end{align*}
پس اگر $f$ و~$g$ بازگشتی اولیه باشند، $\Add$ نیز چنین است.
$f(x_0)=x_0=\Proj{1}{0}(x_0)$ و توابع افکنش بازگشتی اولیه به شمار
می‌روند؛ پس $f$ بازگشتی اولیه است. تابع $g$ تابع سه‌موضعیِ
$g(x_0,y,z)$ است که چنین تعریف می‌شود:
\[
g(x_0, y, z) = \Succ(z).
\]
این هنوز نشان نمی‌دهد که $g$ بازگشتی اولیه است، زیرا $g$ و~$\Succ$
کاملاً یک تابع نیستند: $\Succ$ یک‌موضعی است و $g$ باید سه‌موضعی باشد.
اما می‌توانیم $g$ را به‌طور «رسمی» با ترکیب چنین تعریف کنیم:
\[
g(x_0, y, z) = \Succ(\Proj{3}{2}(x_0, y, z))
\]
چون $\Succ$ و~$\Proj{3}{2}$ توابع بازگشتی اولیه به شمار می‌روند،
$g$ نیز چنین است، زیرا با ترکیب از توابع بازگشتی اولیه تعریف می‌شود.
\end{proof}

\begin{prop}
  \ollabel{prop:mult-pr}
  تابع ضرب $\Mult(x,y)=x\cdot y$ بازگشتی اولیه است.
\end{prop}

\begin{proof}
  تمرین.
\end{proof}

\begin{prob}
  \olref[cmp][rec][prf]{prop:mult-pr} را ثابت کنید: نشان دهید تعریف
  بازگشتی اولیهٔ~$\Mult$ را می‌توان در قالب لازمِ
  \olref[cmp][rec][prf]{defn:primitive-recursion} نوشت و توابع متناظر
  $f$ و~$g$ بازگشتی اولیه‌اند.
\end{prob}

\begin{ex}
این نخستین مثال ما از تعریف با بازگشت اولیه بود:
\begin{align*}
h(0) & =  1 \\
h(y+1) & =  2 \cdot h(y).
\end{align*}
این تابع در قالب لازمِ \olref{defn:primitive-recursion} نمی‌گنجد،
زیرا $k=0$ است. تعریف همچنین ثابت‌های $1$ و~$2$ را دربر دارد. برای
رفع مشکل نخست، آرگومانی صوری می‌افزاییم و تابع~$h'$ را تعریف می‌کنیم:
\begin{align*}
h'(x_0, 0) & =  f(x_0) = 1 \\
h'(x_0, y+1) & =  g(x_0, y, h'(x_0, y)) = 2 \cdot h'(x_0, y).
\end{align*}
تابع $f(x_0)=1$ را می‌توان با ترکیب از $\Succ$ و~$\Zero$ تعریف کرد:
$f(x_0)=\Succ(\Zero(x_0))$. تابع $g$ را می‌توان با ترکیب از
$g'(z)=2\cdot z$ و افکنش‌ها تعریف کرد:
\begin{align*}
g(x_0, y, z) & = g'(\Proj{3}{2}(x_0, y, z))
\intertext{و خودِ $g'$ را می‌توان با ترکیب چنین تعریف کرد:}
g'(z) & = \Mult(g''(z), \Proj{1}{0}(z))
\intertext{و}
g''(z) & = \Succ(f(z)),
\end{align*}
که در آن $f$ همان تابع بالاست: $f(z)=\Succ(\Zero(z))$. اکنون که
$h'$ را داریم، دوباره از ترکیب استفاده می‌کنیم و می‌گذاریم
$h(y)=h'(\Proj{1}{0}(y),\Proj{1}{0}(y))$. این نشان می‌دهد که $h$ را
می‌توان با دنباله‌ای از ترکیب‌ها و بازگشت‌های اولیه از توابع پایه
تعریف کرد؛ پس $h$ بازگشتی اولیه است.
\end{ex}

\end{document}
